% Canonical BandNet diatomic derivation.
% Requires generated/derivations/equations.tex.

This derivation treats a one-dimensional diatomic lattice with local and
beyond-nearest-neighbor interactions, following the nonlocal spring-mass lattice
formulation of Ref.~\cite{kazemi2023drawing}. It proceeds from the site equation
of motion to the two cell equations, assembles the Bloch eigenproblem, solves its
generalized determinant for the acoustic and optical branches, and closes with
consistency checks.

\subsection{Model and notation}

Consider an infinite, free-standing chain of point masses in which two positive
masses $m_1$ and $m_2$ alternate on equally spaced sites $x_j=jd$,
$j\in\mathbb{Z}$; no assumption $m_1\neq m_2$ is made. The repeat period is
$L=2d$: cell $\ell\in\mathbb{Z}$ contains an $m_1$ mass at $x=2\ell d$ and an
$m_2$ mass at $x=(2\ell+1)d$. Every mass is connected by a passive linear spring
of stiffness $k_n\geq0$ to the two masses $n$ sites away,
$n=1,2,\ldots,N$, and no spring connects any mass to ground. Even $n$ connects
sites on the same sublattice; odd $n$ connects opposite sublattices. Let
$u_\ell$ and $v_\ell$ denote the two displacements in cell $\ell$.

\subsection{Equations of motion}

For a site $j$ with displacement $w_j$ and mass $m$, balance of linear momentum
gives
\begin{equation}
    m\ddot{w}_j=\sum_{n=1}^{N}k_n
    \left(w_{j+n}+w_{j-n}-2w_j\right).
    \label{eq:site}
\end{equation}
Separating the even and odd interactions gives
\begin{align}
    m_1\ddot{u}_\ell
    &=\sum_{\substack{n=2\\\text{$n$ even}}}^{N}k_n
      \left(u_{\ell+\frac n2}+u_{\ell-\frac n2}-2u_\ell\right)
      +\sum_{\substack{n=1\\\text{$n$ odd}}}^{N}k_n
      \left(v_{\ell+\frac{n-1}{2}}+v_{\ell-\frac{n+1}{2}}-2u_\ell\right),
      \label{eq:cellA}\\
    m_2\ddot{v}_\ell
    &=\sum_{\substack{n=2\\\text{$n$ even}}}^{N}k_n
      \left(v_{\ell+\frac n2}+v_{\ell-\frac n2}-2v_\ell\right)
      +\sum_{\substack{n=1\\\text{$n$ odd}}}^{N}k_n
      \left(u_{\ell+\frac{n+1}{2}}+u_{\ell-\frac{n-1}{2}}-2v_\ell\right).
      \label{eq:cellB}
\end{align}

\subsection{Bloch ansatz and eigenproblem}

Use the cell-index Bloch gauge
\begin{equation}
    u_\ell(t)=Ue^{i(q\ell L-\omega t)},\qquad
    v_\ell(t)=Ve^{i(q\ell L-\omega t)}.
    \label{eq:bloch}
\end{equation}
The intra-cell offset is absorbed into $V$. A position-based gauge moves the
factors $e^{\pm iqd}$ from the matrix into the amplitude without changing the
eigenfrequencies. Each even interaction contributes
$2[\cos(nqd)-1]$. Each odd interaction contributes
\begin{equation*}
    e^{iq\frac{n-1}{2}L}+e^{-iq\frac{n+1}{2}L}
    =2e^{-iqd}\cos(nqd)
\end{equation*}
in the first equation and its complex conjugate in the second. Therefore
\begin{equation}
    \BandNetDiatomicEigenproblem.
    \label{eq:eigen}
\end{equation}
Here
\begin{equation}
    \BandNetDiatomicKDefinitions.
    \label{eq:dia-k}
\end{equation}
The matrix on the right is Hermitian, so its generalized eigenvalues are real
for positive masses.

\subsection{Dispersion branches}

Nontrivial solutions require
\begin{equation}
    \BandNetDiatomicDeterminant.
    \label{eq:determinant}
\end{equation}
Solving this quadratic in $\omega^2$ gives
\begin{equation}
    \BandNetDiatomicBranches.
    \label{eq:dia-branches}
\end{equation}
The minus and plus signs define the lower and upper branches. When at least one
odd-distance stiffness is positive, these are conventionally the acoustic and
optical branches. Using
$(1/m_1+1/m_2)^2-4/(m_1m_2)=(1/m_1-1/m_2)^2$, the radicand is
\begin{equation}
    \BandNetDiatomicRadicand.
    \label{eq:radicand}
\end{equation}
Thus the branches are real and ordered. Nonnegative stiffnesses also imply
$2K_0\geq|K_1|$, so both squared frequencies are nonnegative.

\subsection{Wave-number domain and consistency checks}

The first Brillouin zone is $|q|\leq\pi/(2d)$. The dispersion is even, so the
positive half is represented by $\hat q=2qd/\pi\in[0,1]$.

\paragraph{Zone center.}
Let $K_{\mathrm{odd}}=\sum_{\text{$n$ odd}}k_n$. At $q=0$,
$K_0=K_{\mathrm{odd}}$ and $K_1=2K_{\mathrm{odd}}$, giving
\begin{equation*}
    \omega_{\mathrm{ac}}^2(0)=0,\qquad
    \omega_{\mathrm{op}}^2(0)=2K_{\mathrm{odd}}
    \left(\frac1{m_1}+\frac1{m_2}\right).
\end{equation*}
The zero acoustic value reflects the absence of grounded springs. The optical
gap is set entirely by odd-$n$ springs between the two sublattices.

\paragraph{Zone edge.}
At $qd=\pi/2$, all odd-$n$ cosine terms vanish. The sublattices decouple and the
two values are $2K_0/m_1$ and $2K_0/m_2$. Their squared-frequency separation is
$2K_0|1/m_1-1/m_2|$ and closes when $m_1=m_2$.

\paragraph{Nearest-neighbor limit.}
If only $k_1$ remains, $K_0=k_1$ and $K_1=2k_1\cos(qd)$. The branches reduce to
\begin{equation*}
    \omega^2(q)=k_1\left(\frac1{m_1}+\frac1{m_2}\right)
    \mathbin{\pm}k_1\sqrt{
        \left(\frac1{m_1}+\frac1{m_2}\right)^2
        -\frac{4\sin^2(qd)}{m_1m_2}},
\end{equation*}
the standard nearest-neighbor diatomic dispersion.

\paragraph{Equal masses.}
For $m_1=m_2=m$, the radicand is $K_1^2/m^2$ and the ordered branches are
$(2K_0\mp|K_1|)/m$. The unordered pair can be written
\begin{equation*}
\left\{
\begin{aligned}
    \frac{2K_0-K_1}{m}
        &=\frac2m\sum_{n=1}^{N}k_n\left[1-\cos(nqd)\right],\\
    \frac{2K_0+K_1}{m}
        &=\frac2m\sum_{n=1}^{N}k_n
          \left[1-\cos\left(n(\pi-qd)\right)\right].
\end{aligned}
\right.
\end{equation*}
Here $\cos(n(\pi-qd))=(-1)^n\cos(nqd)$. These are the monoatomic dispersion at
$q$ and at the folded point $\pi/d-q$. The curves can exchange lower and upper
ordering where $K_1$ changes sign and meet where $K_1=0$, including the zone
edge.
